\begin{thebibliography}{99}
\bibitem{tsai} Tsai Y.-H., Hung W.-C., Schulter S. et al. Learning to Adapt Structured Output Space for Semantic Segmentation // CVPR. 2018. P.~7472--7481. \href{https://openaccess.thecvf.com/content_cvpr_2018/html/Tsai_Learning_to_Adapt_CVPR_2018_paper.html}{Текст публикации}.
\bibitem{dacs} Tranheden W., Olsson V., Pinto J., Svensson L. DACS: Domain Adaptation via Cross-Domain Mixed Sampling // WACV. 2021. P.~1379--1389. \href{https://openaccess.thecvf.com/content/WACV2021/html/Tranheden_DACS_Domain_Adaptation_via_Cross-Domain_Mixed_Sampling_WACV_2021_paper.html}{Текст публикации}.
\bibitem{daformer} Hoyer L., Dai D., Van Gool L. DAFormer: Improving Network Architectures and Training Strategies for Domain-Adaptive Semantic Segmentation // CVPR. 2022. P.~9924--9935. \href{https://openaccess.thecvf.com/content/CVPR2022/html/Hoyer_DAFormer_Improving_Network_Architectures_and_Training_Strategies_for_Domain-Adaptive_Semantic_CVPR_2022_paper.html}{Текст публикации}.
\bibitem{askari} Аскари Хеммат С. Применение синтетических данных из UnrealEngine для обучения модели сегментации мебели // Искусственный интеллект в промышленных, коммерческих, медицинских и финансовых приложениях. Вып.~4. М.: НИТУ «МИСИС», 2025. С.~13--16. \href{https://sadekov.su/Articles/kik_2025_1.pdf#page=14}{Сборник, с.~13}.
\bibitem{komendala} Комендала Л.~Б. Семантическая сегментация дорожных сцен с использованием глубокого обучения // Искусственный интеллект в промышленных, коммерческих, медицинских и финансовых приложениях. Вып.~5. М.: НИТУ «МИСИС», 2026. С.~82--86. \href{https://sadekov.su/Articles/kik_2025_2.pdf#page=83}{Сборник, с.~82}.
\bibitem{kurochkin} Kurochkin I.~I., Kurochkin I.~N., Kolosova O.~Yu., Lozinsky V.~I. Cryostructuring of Polymeric Systems: Application of Deep Neural Networks for the Classification of Structural Features Peculiar to Macroporous Poly(vinyl alcohol) Cryogels Prepared without and with the Additives of Chaotropes or Kosmotropes // Molecules. 2020. Vol.~25, no.~19. Art.~4480. \href{https://doi.org/10.3390/molecules25194480}{DOI: 10.3390/molecules25194480}.
\bibitem{mmd} Gretton A., Borgwardt K.~M., Rasch M.~J. et al. A Kernel Two-Sample Test // JMLR. 2012. Vol.~13. P.~723--773. \href{https://www.jmlr.org/papers/v13/gretton12a.html}{Текст публикации}.
\bibitem{coral} Sun B., Saenko K. Deep CORAL: Correlation Alignment for Deep Domain Adaptation. 2016. \href{https://arxiv.org/abs/1607.01719}{arXiv:1607.01719}.
\bibitem{sinkhorn} Cuturi M. Sinkhorn Distances: Lightspeed Computation of Optimal Transport // NeurIPS. 2013. Vol.~26. \href{https://arxiv.org/abs/1306.0895}{arXiv:1306.0895}.
\bibitem{advent} Vu T.-H., Jain H., Bucher M. et al. ADVENT: Adversarial Entropy Minimization for Domain Adaptation in Semantic Segmentation // CVPR. 2019. \href{https://arxiv.org/abs/1811.12833}{arXiv:1811.12833}.
\bibitem{mic} Hoyer L., Dai D., Wang H., Van Gool L. MIC: Masked Image Consistency for Context-Enhanced Domain Adaptation // CVPR. 2023. P.~11721--11732. \href{https://arxiv.org/abs/2212.01322}{arXiv:2212.01322}.
\bibitem{sepico} Xie B., Li S., Li M. et al. SePiCo: Semantic-Guided Pixel Contrast for Domain Adaptive Semantic Segmentation // IEEE TPAMI. 2023. \href{https://arxiv.org/abs/2204.08808}{arXiv:2204.08808}.
\bibitem{gta} Richter S.~R., Vineet V., Roth S., Koltun V. Playing for Data: Ground Truth from Computer Games // ECCV. 2016. \href{https://arxiv.org/abs/1608.02192}{arXiv:1608.02192}.
\bibitem{cityscapes} Cordts M., Omran M., Ramos S. et al. The Cityscapes Dataset for Semantic Urban Scene Understanding // CVPR. 2016. P.~3213--3223. \href{https://www.cityscapes-dataset.com/}{Сайт набора данных}.
\end{thebibliography}
